\documentclass[11pt]{article}
\usepackage[a4paper,margin=1in]{geometry}
\usepackage{fontspec}
\usepackage[boldfont=false]{xeCJK}
\setCJKmainfont{FandolSong-Regular.otf}
\usepackage{amsmath}
\usepackage{amssymb}
\usepackage{array}
\usepackage{booktabs}
\usepackage{graphicx}
\usepackage{adjustbox}
\usepackage{enumitem}
\setlist[itemize]{topsep=2pt,partopsep=0pt,parsep=0pt,itemsep=2pt,leftmargin=1.4em}
\setlist[enumerate]{topsep=2pt,partopsep=0pt,parsep=0pt,itemsep=2pt,leftmargin=1.6em}
\usepackage{parskip}
\usepackage{fvextra}
\fvset{breaklines=true,breakanywhere=true,fontsize=\small}
\setlength{\parindent}{0pt}

\begin{document}

\begin{center}
  {\LARGE\bfseries Probability}\\[0.35em]
  {\large Igcse Mathematics}
\end{center}
\vspace{0.6em}

This handout covers Topic 8, Probability. Parts marked \textbf{(Extended)} are only tested on the Extended papers; everything else is for both levels.

\section*{The probability scale}

\textbf{Probability} 概率 measures how likely an \textbf{event} 事件 is. It runs on a scale from $0$ (impossible) to $1$ (certain), and can be written as a fraction, decimal or percentage. We write $\text{P}(A)$ for the probability of event $A$.

For equally likely \textbf{outcomes} 结果,

\[
\text{P}(\text{event}) = \frac{\text{number of favourable outcomes}}{\text{total number of outcomes}}.
\]

The probability that an event does \textbf{not} happen is

\[
\text{P}(A') = 1 - \text{P}(A).
\]

\textbf{Worked example.} A bag has $3$ red and $5$ blue counters. Find the probability of \emph{not} drawing red.

\[
\text{P}(\text{red}) = \frac{3}{8}, \qquad \text{P}(\text{not red}) = 1 - \frac{3}{8} = \frac{5}{8}.
\]

\section*{Relative frequency and expected frequency}

When outcomes are not equally likely, do an experiment. The \textbf{relative frequency} 相对频率 estimates the probability:

\[
\text{relative frequency} = \frac{\text{number of times it happened}}{\text{total number of trials}}.
\]

The more trials you do, the better the estimate. A \textbf{fair} 公平 object gives equal chances; one with \textbf{bias} 偏倚 does not; \textbf{random} 随机 means each outcome happens by chance.

The \textbf{expected frequency} 期望频数 is how many times you expect an event in $n$ trials:

\[
\text{expected frequency} = \text{P}(\text{event}) \times n.
\]

\textbf{Worked example.} The probability of rolling a six is $\tfrac{1}{6}$. How many sixes are expected in $300$ rolls?

\[
\frac{1}{6} \times 300 = 50.
\]

\section*{Combined events}

For two or more events together (\textbf{combined events} 组合事件), two rules help:

\begin{itemize}
\item \textbf{AND} (both happen): \textbf{multiply} the probabilities — when the events are \textbf{independent} 独立事件 (one does not affect the other).

\item \textbf{OR} (either happens): \textbf{add} the probabilities — when the events are \textbf{mutually exclusive} 互斥 (they cannot both happen).

\end{itemize}

Three pictures help you organise the work.

\subsection*{Sample space diagrams}

A \textbf{sample space diagram} 样本空间图 is a table or grid listing every possible outcome — useful for two dice or two spinners. Count the outcomes you want out of the total.

\subsection*{Venn diagrams}

A \textbf{Venn diagram} 维恩图 sorts outcomes into overlapping sets. From it you can read $\text{P}(A \cap B)$ (in both) and $\text{P}(A \cup B)$ (in either).

\subsection*{Tree diagrams}

A \textbf{tree diagram} 树状图 shows each stage as a set of branches. Write the probability on each branch and the outcome at the end. \textbf{Multiply along} the branches, then \textbf{add} the paths you want.

\textbf{Worked example (with replacement).} From the bag ($3$ red, $5$ blue), a counter is drawn, replaced, then a second is drawn. This is \textbf{with replacement} 有放回, so the chances do not change. Find the probability of two reds.

\[
\text{P}(\text{red, red}) = \frac{3}{8} \times \frac{3}{8} = \frac{9}{64}.
\]

\textbf{Worked example (without replacement, Extended).} Now the first counter is \emph{not} replaced — drawing \textbf{without replacement} 无放回. After one red is taken, $2$ reds remain out of $7$:

\[
\text{P}(\text{red, red}) = \frac{3}{8} \times \frac{2}{7} = \frac{6}{56} = \frac{3}{28}.
\]

\section*{Conditional probability (Extended)}

\textbf{Conditional probability} 条件概率 is the probability of one event \emph{given that} another has already happened. Read it from a Venn diagram, two-way table or tree diagram by looking only at the part that matches the condition.

\textbf{Worked example.} In a class of $30$, $18$ study French and, of those, $7$ also study German. A French student is picked. Find the probability they also study German.

Look only at the $18$ French students:

\[
\text{P}(\text{German} \mid \text{French}) = \frac{7}{18}.
\]


\end{document}
